“Want teveel frustratie
leidt vaak tot rebellie.

En teveel rebellie
leidt vaak tot confrontatie.

Teveel satisfactie
leidt tot resignatie.

Teveel resignatie
leidt tot deprivatie.

Teveel arbeidsschaarste
leidt het kostbaarste

naar onbereikbaarheid.
Een heilige waarheid.”

De administrateur
sprak als een adviseur

op zoek naar een balans
tussen de stimulans

en de relaxatie
binnen de parochie.

Winst maximaliseren
vroeg om te laveren.

Schaarste kweekte welvaart.
Schaarste kweekte kruisvaart.

Cijfers deden ertoe.
De stad was zijn melkkoe.

Hij moest in kunnen staan
voor de stads voortbestaan.
Bij zeurende ergernis;
“Ik zeg niet wat waar is,

ik zeg slechts wat er klopt.”
Kritiek werd afgestopt.

Na een half jaar Emmy,
wist hij haar prestatie.

Op zondag waarachtig
maar liefst honderdtachtig.

Zondag was geen rustdag.
Dat was wel zaterdag.

De andere dagen
moest ze 60 dragen.

Gemiddeld uiteraard.
“Emmy’s misbruik standaard

meet ik elke week weer.
Vierhonderd-tachtig keer

bepaal ik als plafond.”
Dit was de quaestor vond.

Gelaakt door anderen.
Niets zou veranderen

omdat de kerk goed wist
dat Emmy dan beslist

schade zou oplopen
als ze vaker open

voor paring moest gaan staan.
De cijfers toonden aan

dat men medehoeren,
minder vaak kon vloeren.

Door het hoge verloop
had men in de verkoop,

meestal maar twee hoeren
om het uit te voeren.

Geveld door uitputting
of de constatering

van een geslachtsziekte.
Wat direct kortwiekte

de vleugels van plezier,
volgens de thesaurier.

“Het probleem is de wet.
Menig vrouw moet naar bed

met iedere jonkman,
die tevens een buurman

of overbuurman is.
En daar gaat het vaak mis.

Die vrouwenverslinder
zorgt dat er veel minder

met de eigen gemaal
wordt geneukt dan normaal.

Uit de gegevens blijkt
dat vaak verongelijkt

de man elders troost zoekt.
En dan een hoer bezoekt.

Echter vijftig procent
heeft dit probleem herkend

en gaat maar verhuizen,
zodat men naast huizen

met slechts gehuwden woont.
Als de kerk dit beloond

met financiële hulp
dan houdt de man zijn gulp

buitenshuis vaker dicht.”

Dat werkt dubbelop.
De man mocht er op

zo vaak als hij wilde
en dat amper stilde

bij de seksslavinnen.
Voor zich winnen

door verliefd te maken
of een vrouw te schaken

buiten de stadsmuren,
zo vulden de buren

het vrouwentekort aan.
Daar was behoefte aan.

Als 800 mannen
met snode plannen

ontevreden zijn
over hun welzijn

dan wordt het gevaarlijk.
Dat was zo bezwaarlijk

dat de onoprechter
een voorvechter

werd zelfs trouwe
gehuwde vrouwen

vals te veroordelen
om ze te herindelen

bij de parochie
naast Emmy.
